% Generated draft: requires mathematical and source review.
\documentclass[11pt]{article}
\usepackage[margin=25mm]{geometry}
\usepackage{fontspec}
\setmainfont{Latin Modern Roman}
\usepackage{amsmath,amssymb,mathtools,unicode-math}
\setmathfont{Latin Modern Math}
\usepackage{enumitem,needspace,xurl,hyperref}
\setlength{\parindent}{0pt}
\setlength{\parskip}{5pt}
\setlength{\emergencystretch}{3em}
\newcommand{\N}{\mathbb{N}}
\newcommand{\Z}{\mathbb{Z}}
\newcommand{\Q}{\mathbb{Q}}
\newcommand{\R}{\mathbb{R}}
\newcommand{\C}{\mathbb{C}}
\newcommand{\F}{\mathbb{F}}
\newcommand{\A}{\mathbb{A}}
\newcommand{\PP}{\mathbb{P}}
\newcommand{\E}{\mathbb{E}}
\newcommand{\eps}{\varepsilon}
\newcommand{\dd}{\,\mathrm{d}}
\begin{document}
\section*{UnsolvedMath ID: 1211}
\textbf{Source problem code:} catalogue entry
\subsection{Definitions and mathematical statement}
Let $(M,g)$ be a smooth Riemannian manifold of dimension $n$, and let $R$ be its Riemann curvature tensor, with any fixed sign convention. For $p\in M$ and a unit vector $v\in T_pM$, the Jacobi operator on the orthogonal complement of $v$ is
\[
 J_v:v^\perp\longrightarrow v^\perp,\qquad J_v(w)=R(w,v)v.
\]
It is self-adjoint. Write $\operatorname{Spec}(J_v)$ for its multiset of eigenvalues, including multiplicities. Using $T_pM$ instead of $v^\perp$ merely adjoins the always-present eigenvalue corresponding to $v$ and does not change the invariance condition.

The manifold is \emph{pointwise Osserman} if
\[
 \forall p\in M\ \exists\Lambda_p\ \forall v\in T_pM:
 \quad g_p(v,v)=1\Longrightarrow\operatorname{Spec}(J_v)=\Lambda_p.
\]
It is \emph{globally Osserman} if there is a single multiset independent of both point and unit direction:
\[
 \exists\Lambda\ \forall p\in M\ \forall v\in T_pM:
 \quad g_p(v,v)=1\Longrightarrow\operatorname{Spec}(J_v)=\Lambda.
\]
Thus pointwise constancy need not include constancy as $p$ varies. The supplied background gives the pointwise condition, while the dated review identifies the classical conjecture as using the global condition.

A manifold is flat if $R$ vanishes identically. A rank-one symmetric space here means a nonflat irreducible Riemannian symmetric space whose maximal totally geodesic flat subspaces have dimension one, with its metric scale included as part of the model. Being locally isometric to such a space means every point has a neighborhood isometric to an open subset of that model; no global isometry, completeness, or simple-connectedness assumption is thereby imposed.

\textbf{Imported question.} Is every Osserman Riemannian manifold either flat or locally isometric to a rank-one symmetric space? With the background's definition, this asks the question for pointwise Osserman manifolds. The corresponding \emph{global} Osserman conjecture asks the same conclusion under the stronger, globally constant spectrum condition. These formulations must be distinguished; this draft does not replace the background definition without notice or infer an unrestricted pointwise theorem from global results. Connectedness is not specified in the record; whether conclusions are to be interpreted componentwise or for a single connected manifold must also be fixed in a final formulation.

The literature triage dated 2026-08-17 reports that Nikolayevsky proved the global conjecture outside dimensions $8$ and $16$, and separately settled dimension $8$, leaving dimension $16$ as the exceptional unresolved global case in that assessment. It says the older background's claim of knowledge only through dimension $4$ is obsolete, and calls for distinguishing pointwise and global definitions. The cited work is said to state further pointwise qualifications, but those qualifications are not reproduced in the supplied record. Consequently they are not reconstructed here. The registry's \emph{open} label and the triage's \emph{partially solved} label are supplied dated descriptions, not a fresh literature verification.
\subsection{Short English statement}
The Osserman conjecture asks whether constancy of the Jacobi-operator eigenvalues forces a Riemannian manifold to be flat or locally rank-one symmetric. The record conflates constancy at each point with constancy over all points. Its dated review reports that the global conjecture is settled except in dimension $16$, without specifying all qualifications for the pointwise version.
\subsection{Sources}
\begin{itemize}
\item UnsolvedMath local registry, data/unsolvedmath/problems.json, permanent ID 1211, source code catalogue entry, ``The Osserman Conjecture.'' The background attributes the conjecture to Robert Osserman and supplies the pointwise definition; the literature triage is dated 2026-08-17.
\item Yuri Nikolayevsky, ``Osserman Conjecture in dimension $n\neq8,16$,'' \emph{Mathematische Annalen} 331 (2005), 505--522, supplied link \url{https://arxiv.org/abs/math/0204258}. The review attributes the global result outside dimensions $8,16$ and corresponding pointwise qualifications to this paper.
\item Yuri Nikolayevsky, ``Osserman manifolds of dimension 8,'' \emph{Manuscripta Mathematica} 115 (2004), 31--53, supplied link \url{https://arxiv.org/abs/math/0310387}. The review cites this as settling dimension $8$.
\item Yuri Nikolayevsky, ``Conformally Osserman manifolds,'' \emph{Pacific Journal of Mathematics} 245 (2010), 315--337, \url{https://doi.org/10.2140/pjm.2010.245.315}. The review cites its statement that dimension $16$ remains exceptional for the original conjecture.
\item Ulam AI and the UnsolvedMath Contributors, \emph{UnsolvedMath: A Curated Collection of Open Mathematics Problems}. Dataset curation and original metadata are supplied under CC BY 4.0; referenced source material retains its own terms. No linked reference was independently consulted for this draft.
\end{itemize}
\end{document}
